%============================================================
%  Bd. 25 - Zusammenhängende Graphen
%============================================================

\documentclass[main.tex]{subfiles}

\ifSubfilesClassLoaded{
  \BandSetupStandalone{B25}
}{
}

\title{Bd. 25 - Zusammenhängende Graphen}
\author{Martin Kunze}
\date{}
\setcounter{file}{25}

\hypersetup{
  pdftitle={Bd. 25 - Zusammenhängende Graphen},
  pdfauthor={Martin Kunze}
}

\begin{document}

\title{Bd. 25 - Zusammenhängende Graphen}
\author{Martin Kunze}
\date{}
\hypersetup{pageanchor=false}
\maketitle
\hypersetup{pageanchor=true}
\setcounter{tocdepth}{3}
\tableofcontents
\setcounter{file}{25}
\renewcommand{\FormulaBandID}{25}

\chapter{Einleitung}

Ein ungerichteter Graph heißt zusammenhängend, wenn je zwei seiner Knoten
durch einen endlichen Walk verbunden sind. Der in Band~22 eingeführte
Erreichbarkeitsbegriff enthält dabei auch den Fall der Länge null und erfasst
somit den Fall gleicher Endpunkte ohne Sonderregel.

Dieser Band isoliert den Zusammenhang als eigene Strukturstufe. Dadurch
kann die folgende Baumtheorie die Zusammenhangseigenschaft unabhängig von
der in Band~24 eingeführten Schleifenfreiheit verwenden.

\chapter{Zusammenhängende Graphen}

\section{Axiomatischer Begriff}

\FormulaDefDeltaK[Begriff des zusammenhängenden Graphen]{%
  \ConnectedGraphStruct{V,E}%
}{ConnectedGraphDef}{%
  \DeltaRow{Mengen}{V\dsep E\dsep u\dsep v}
  \DeltaRow{\textbf{Axiome}}{}
  \DeltaRow{Ungerichteter Graph}{\UndirectedGraphStruct{V,E}}
    [\FormulaRefAuto{UndirectedGraphDef}[def]]
  \DeltaRow{Erreichbarkeit}{%
    u,v\in V\vdash u\mathrel{\GraphReach{E}}v}
    [\FormulaRefAuto{DirectedReachabilityDef}[def]]
  \DeltaRow{\textbf{Neues Symbol}}{}
  \DeltaPrem{Zusammenhängender Graph}{\ConnectedGraphStruct{V,E}}
}

\FormulaAxiomDeltaK[Zugriff auf den ungerichteten Grundgraphen]{%
  \ConnectedGraphStruct{V,E}\vdash\UndirectedGraphStruct{V,E}%
}{ConnectedGraphUnderlyingAxiom}{%
  \DeltaRow{Mengen}{V\dsep E}%
}

\FormulaAxiomDeltaK[Zugriff auf den Zusammenhang]{%
  \ConnectedGraphStruct{V,E}\vdash
  \forall u,v\in V\,\bigl(u\mathrel{\GraphReach{E}}v\bigr)%
}{ConnectedGraphReachabilityAxiom}{%
  \DeltaRow{Mengen}{V\dsep E\dsep u\dsep v}%
}

\FormulaThmDeltaK[Strukturzugriff auf einen zusammenhängenden Graphen]{%
  \ConnectedGraphStruct{V,E}\vdash
  \UndirectedGraphStruct{V,E}\land\GraphStruct{V,E}%
}{ConnectedGraphAccess}{%
  \DeltaRow{Mengen}{V\dsep E}%
}
\begin{tabproof}
  \proofstep{1}{\ConnectedGraphStruct{V,E}}{\rA}
  \proofstep{1}{\UndirectedGraphStruct{V,E}}{%
    \FormulaRefAuto{ConnectedGraphUnderlyingAxiom}[ax]{1}}
  \proofstep{1}{\GraphStruct{V,E}}{%
    \FormulaRefAuto{UndirectedGraphAccess}[thm]{2}}
  \proofstep{1}{\UndirectedGraphStruct{V,E}\land\GraphStruct{V,E}}{\rAI{2,3}}
\end{tabproof}

\FormulaThmDeltaK[Je zwei Knoten eines zusammenhängenden Graphen sind erreichbar]{%
  \ConnectedGraphStruct{V,E}\dsep u,v\in V\vdash
  u\mathrel{\GraphReach{E}}v%
}{ConnectedGraphVertexReachability}{%
  \DeltaRow{Mengen}{V\dsep E\dsep u\dsep v}%
}
\begin{tabproof}
  \proofstep{1}{\ConnectedGraphStruct{V,E}}{\rA}
  \proofstep{2}{u\in V}{\rA}
  \proofstep{3}{v\in V}{\rA}
  \proofstep{1}{%
    \forall u,v\in V\,\bigl(u\mathrel{\GraphReach{E}}v\bigr)}{%
    \FormulaRefAuto{ConnectedGraphReachabilityAxiom}[ax]{1}}
  \proofstep{1,2,3}{u\mathrel{\GraphReach{E}}v}{%
    \rRE{3,\rUE{\rRE{2,\rUE{4}}}}}
\end{tabproof}

\end{document}
